% !TEX root = main07_en.tex
\documentclass[reprint,amsmath,amssymb,aps,prd]{revtex4-2}

\usepackage{graphicx}
\usepackage{dcolumn}
\usepackage{bm}
\usepackage{hyperref}

\begin{document}

\preprint{APS/123-QED}

\title{Mechanics of Spatial Vibration III: Cosmological Spatial Plate Tectonics and Wave Interference Topology}

\author{Kwang Yong Yoo}
\email{khanyong.yoo@gmail.com}
\affiliation{Independent Researcher \\ 
\textit{AI-Assisted Research and Simulation Development Support: Gemini} \\
(Interactive simulations and supplementary materials are officially archived under DOI: \href{https://doi.org/10.5281/zenodo.21258029}{10.5281/zenodo.21258029})}

\date{July 8, 2026}

\begin{abstract}
As the concluding part of the Spatial Vibration Mechanics series, this paper addresses the inhomogeneity of the universe by introducing the speculative hypothesis of \textbf{`Cosmological Spatial Plate Tectonics'}.
Due to the causality limits of its accelerated expansion, the universe is hypothesized to be fragmented into colossal `Topological Domains' possessing distinct intrinsic vibrational phases.
By analyzing the wave interference mechanics at the `Spatial Fault Lines' where these domains converge, this study offers a phenomenological geometric interpretation for the formation of Cosmic Voids and the Cosmic Web.
To preserve local energy-momentum conservation, destructive interference is modeled not as energy annihilation, but as the hydrodynamic redistribution of tensor energy.
Furthermore, we model \textbf{`Cosmic Quakes'}---ruptures caused by topological friction stress tensors ($S_{\mu\nu}$) exceeding the elastic limit of spacetime.
We propose that these catastrophic tensor ruptures emit high-frequency gravitational bursts, which convert into Fast Radio Bursts (FRBs) via the Gertsenshtein effect.
In extreme cases, this rupture leads to macroscopic topological pinch-off, resulting in the spatial separation of macroscopic domains and the continuous genesis of multiverses.
Finally, utilizing the anisotropy of the Stochastic Gravitational-Wave Background (SGWB), this study proposes the \textbf{POINTING} protocol framework to map the 3D topology of cosmic fault lines, establishing a predictive observational methodology for cosmic extreme events.
\end{abstract}

\maketitle

\section{Introduction}
Previous papers in this series modeled quantum and macroscopic dark phenomena through the `geometrical fluctuation of space'.
However, applying this model to the entire expanding universe confronts a fundamental phase-synchronization challenge.
The rapid expansion of the universe \cite{Guth1981} implies that distant local regions are causally disconnected.
Therefore, assuming the entire universe maintains a uniform phase coherence of spatial vibration is physically problematic.
To resolve this, this study hypothesizes that the expanding cosmic space has fragmented into 3D blocks---\textbf{`Topological Domains'}---which maintain internal coherence but differ in vibrational phase from adjacent regions.
This concept is termed \textbf{`Cosmological Spatial Plate Tectonics'}.

\section{Cosmological Spatial Plate Tectonics}
This section phenomenologically formulates the interaction mechanics of spatial domains with differing phases.

\subsection{Spatial Plates and Fault Lines}
An independent spatial region where the vibration phase angle ($\theta_i$) is synchronized is defined as a Spatial Plate.
At macroscopic boundaries where plates with different phases ($\theta_1, \theta_2$) converge, spacetime curvature experiences topological friction.
This boundary region is termed the \textbf{`Spatial Fault Line ($\Sigma_{1,2}$)'}.
To preserve exact dimensional consistency and the rigorous tensor grammar of general relativity, we introduce a macroscopic coherence length constant ($L_c$).
The effective geometrical stress tensor $S_{\mu\nu}$ accumulated due to the phase mismatch ($\Delta\theta$) is rigorously modeled as:
\begin{equation}
\begin{split}
S_{\mu\nu} &\propto \rho_{\mathrm{vib}} c^2 L_c^2 \bigg( \nabla_\mu(\Delta\theta) \nabla_\nu(\Delta\theta) \\
&\quad - \frac{1}{2} g_{\mu\nu} |\nabla(\Delta\theta)|^2 \bigg)
\end{split}
\end{equation}

\subsection{Resolution of the Domain Wall Problem}
If these fault lines are analogous to topological domain walls, their existence must not severely contradict the observed isotropy of the Cosmic Microwave Background (CMB).
We propose an \textit{Early Dilution Hypothesis}: during the inflationary epoch, the macroscopic energy density of these fault lines was exponentially diluted.
Currently, they exist merely as faint residual stress tensors, acting only as geometrical attractors for dark matter and baryonic gas, safely evading current CMB observational limits.

\section{Topology of Wave Interference and Large-Scale Structures}
This study suggests that the cosmic large-scale structure may be interpreted as a deterministic topology sculpted by the wave interference of spatial plates.

\subsection{Destructive Interference and Energy Redistribution: Cosmic Voids}
Where plates collide with opposing phases ($\Delta\theta \approx \pi$), the spatial vibration undergoes destructive interference.
To strictly preserve local energy-momentum conservation ($\nabla_\mu T^{\mu\nu} = 0$), tensor energy does not simply vanish.
Instead, governed by the covariant 4-dimensional continuity equation, the resulting geometrical pressure gradients force the tensor fluid to be hydrodynamically advected away from the interference node:
\begin{equation}
\lim_{\Delta\theta \to \pi} \left( \tilde{V}_{\mu\nu}^{(1)} + \tilde{V}_{\mu\nu}^{(2)} \right) \implies \nabla_\lambda (\rho_{\mathrm{vib}} u^\lambda_{\mathrm{drift}}) = 0
\end{equation}
As wave cancellation causes geometric pressure to diverge, the tensor fluid is swept outward along its 4-velocity $u^\lambda_{\mathrm{drift}}$ towards the surrounding constructive interference regions (filaments).
In this evacuated region, tensor condensation cannot occur, suppressing matter aggregation.
This provides a formal hydrodynamic origin for vast structures like the Bo\"{o}tes Void \cite{Kirshner1981}.

\subsection{Constructive Interference: Cosmic Web Filaments}
Along linear fault lines where phases nearly match ($\Delta\theta \approx 0$), constructive interference amplifies the vibration energy, emerging as effective surplus gravity. Galaxies accumulate in these fault valleys, forming galactic filaments \cite{Bond1996}. Large-scale attractors could potentially be interpreted as spatial plate nodes where multiple constructive interferences converge.

\section{Cosmic Seismology and Macroscopic Spatial Separation}
Fast Radio Bursts (FRBs) \cite{Lorimer2007} are proposed here to be potentially associated with \textbf{`Cosmic Quakes'} caused by the rupture of spatial fault lines.

\subsection{Phase Locking (Foreshocks) and Fault Rupture (Main Quake)}
Geometric discontinuity locks spacetime tightly at fault boundaries (`Phase Locking'), accumulating elastic potential stress ($S_{\mu\nu}$).
While the fault lines are locked and frictionally grind against each other, this prolonged topological friction emits ultra-low-frequency (nHz band) gravitational noise.
These act as cosmic `foreshocks' that manifest as the Stochastic Gravitational-Wave Background (SGWB).
Conversely, when this accumulated stress exceeds the spacetime elastic critical limit ($S_{\mathrm{crit}}$), a full rupture (`main quake') occurs.
The pent-up tensor energy erupts explosively, dissipating instantaneously as an ultra-high-frequency (GHz band) gravitational wave burst.

\subsection{Electromagnetic Conversion via the Gertsenshtein Effect}
How do high-frequency gravitational ruptures manifest as electromagnetic photons (FRBs) to our telescopes?
We employ the Gertsenshtein effect \cite{Gertsenshtein1961}. The extreme high-frequency (GHz) gravitational wave burst emitted during a cosmic main quake resonates with the vast, strong intergalactic magnetic fields present in surrounding plasmas.
This effectively converts gravitational wave energy into an electromagnetic shockwave (photons), perfectly accounting for the nature of FRBs.
The fault slip-and-relock mechanism naturally explains Repeating FRBs \cite{CHIME2020}.

\subsection{Extreme Rupture: Macroscopic Topological Pinch-off and Spatial Separation}
Expanding the thought experiment to its extreme limits: if the expansive vibrational radiation pressure ($p_{\mathrm{vib}} < 0$) violently tears at the plates, or if the friction stress exponentially exceeds the macroscopic elastic limit, a simple slip (quake) is insufficient.
This induces a \textbf{Macroscopic Topological Pinch-off}.

Classical general relativity forbids macroscopic topology change without singularities (e.g., Geroch's theorem).
Therefore, this macroscopic topological pinch-off implies a phase transition at the Planck scale, where the extreme tensor stress surpasses the elastic limits of the classical metric.
At this threshold, standard spacetime breaks down, spawning a \textbf{Topological Black Hole} that serves as a closed geometric umbilical cord to the isolated baby universe, exactly analogous to the microscopic vortex regularization introduced in Part I.

The severed portion of the spatial plate completely loses its geometric connection to our parent universe, resulting in a profound \textbf{`Spatial Separation'}.
This isolated, massive topological bubble, laden with immense tensor energy, begins independent expansion, phenomenologically serving as a continuous womb for the genesis of a \textbf{Multiverse}.

\section{The POINTING Predictive Protocol}
To observationally falsify and demonstrate this framework, we introduce the \textbf{POINTING} protocol (\textbf{P}hasic \textbf{O}bservation \& \textbf{I}nterferometric \textbf{N}etwork for \textbf{T}ensor-field \textbf{I}nduced \textbf{N}eural \textbf{G}ravitational-signals).
This protocol operates through four distinct pillars: (1) \textit{Phasic Observation} captures the phase-locking phenomena between the SGWB and FRBs;
(2) \textit{Interferometric Network} leverages pulsar timing array (PTA) networks \cite{Agazie2023} for an interferometric approach;
(3) \textit{Tensor-field Induced} addresses the core mechanism of spacetime deformation driven by the spatial vibration model;
and (4) \textit{Neural Gravitational-signals} utilizes machine learning (neural network) pattern recognition to map the non-linear correlations between ultra-low-frequency (nHz) foreshocks and impending high-frequency (GHz) main quakes.

To satisfy rigorous mathematical validity, the rupture probability function ($P_{\mathrm{quake}}$) at fault intersecting nodes is formalized using the scalar norm of the stress tensor ($\| S_{\mu\nu} \|$):
\begin{equation}
\begin{split}
P_{\mathrm{quake}}(x,y,z, t) &\propto \int_{0}^{t} \exp\left( \frac{\| S_{\mu\nu} \| - S_{\mathrm{crit}}}{\alpha} \right) dt' \\
&\quad \left( \text{where } \| S_{\mu\nu} \| \equiv \sqrt{S_{\mu\nu} S^{\mu\nu}} \right)
\end{split}
\end{equation}

By actively monitoring the nHz SGWB foreshocks via the \textbf{POINTING} protocol, radio telescopes can be proactively directed to specific coordinates before an imminent extreme explosion (GHz FRB main quake) occurs, establishing a predictive observational methodology.

\section{Conclusion of the Trilogy}
The Mechanics of Spatial Vibration series presents a unified conceptual framework attributing quantum wave-particle duality, dark sector phenomena, and cosmic large-scale structures to the geometrical fluctuation of space and its scale-dependent interactions with mass inertia.
By interpreting the interference pattern on the microscopic double-slit screen and the cosmic web filaments as physically analogous mechanical phenomena, this model proposes a deterministic foundation for physical reality governed by \textbf{`Fractal Scale Symmetry'}.
The universe does not rely on unknown ghost particles or magic;
particles and galaxies merely navigate along the curvatures of a deterministic symphony played by spacetime tensors.
As a highly speculative framework, establishing the physical validity of Spatial Plate Tectonics requires rigorous relativistic derivations, cosmological N-body simulations incorporating tensor friction, and direct verification through the proposed POINTING protocol.

\vspace{1em}
\section*{ACKNOWLEDGMENTS}
\textbf{Note on Authorship:} This work was authored by Kwang Yong Yoo. The author acknowledges the use of AI-assisted tools (Gemini) for language refinement and simulation development support. This study is an independent research work conducted under the author's sole academic responsibility.

\begin{thebibliography}{99}
\bibitem{Guth1981} A. H. Guth, \textit{Phys. Rev. D} \textbf{23}, 347 (1981).
\bibitem{Kirshner1981} R. P. Kirshner \textit{et al.}, \textit{Astrophys. J.} \textbf{248}, L57 (1981).
\bibitem{Bond1996} J. R. Bond \textit{et al.}, \textit{Nature} \textbf{380}, 603 (1996).
\bibitem{Lorimer2007} D. R. Lorimer \textit{et al.}, \textit{Science} \textbf{318}, 777 (2007).
\bibitem{CHIME2020} CHIME/FRB Collaboration, \textit{Nature} \textbf{582}, 351 (2020).
\bibitem{Agazie2023} G. Agazie \textit{et al.} (NANOGrav Collaboration), \textit{Astrophys. J. Lett.} \textbf{951}, L8 (2023).
\bibitem{Gertsenshtein1961} M. E. Gertsenshtein, \textit{Sov. Phys. JETP} \textbf{14}, 84 (1961).
\end{thebibliography}

\end{document}